
\chapter{Эмитент и эквайер: две стороны одной транзакции}\label{ch:issuer-acquirer}

Когда покупатель прикладывает карту к терминалу или нажимает
кнопку оплаты в интернете, снаружи виден один короткий жест.
Внутри же запускаются два разных процесса. Для эмитента это
мгновенная проверка доверия: действительна ли карта, есть ли деньги,
не слишком ли велик риск мошенничества. Для эквайера это
отдельная задача: принять запрос, провести его по нужному
маршруту и затем корректно рассчитаться с мерчантом.

\section{Эмитент}\label{sec:ia-issuer}

Эмитент (issuer) выпускает карту и принимает на себя кредитный
риск транзакции. В момент покупки именно он должен быстро решить,
готов ли он оплатить операцию через карточную сеть, а затем
взыскать средства с держателя карты, списать их с дебетового
счёта или увеличить задолженность по кредитному. Чтобы такой ответ
занимал секунды, а не дни, у эмитента есть набор систем, каждая
из которых держит свой слой реальности: состояние карты, правила
авторизации, аутентификацию, токены и претензии после операции.

\subsection*{Система управления картой (CMS)}

Система управления картой
(Card Management System, CMS) -- центральная система, управляющая её
жизненным циклом: выпуск, смена статуса, блокировка, перевыпуск. CMS хранит реквизиты
карты и её связь со счётом, поддерживает статусы и лимиты,
формирует задания на персонализацию и фиксирует события,
которые потом будут важны для поддержки, антифрод-мониторинга
и диспутов.

На этапе персонализации на карту загружаются EMV-апплет, ключи эмитента (IMK),
параметры управления риском, такие как floor limit
и офлайн-лимиты, а также данные платёжного приложения
(AID, Application Label) (см. главу~\ref{ch:emv})~\cite{emvco-book3,emvco-book4}. Для карт Мир
персонализация дополнительно включает установку МРА
(Mir Payment Application) с российской криптографией
на базе ГОСТ (см. главу~\ref{ch:mir}).

\subsection*{Авторизационный хост}

Авторизационный хост решает, можно ли по карте провести конкретную
операцию. Когда авторизационный запрос доходит от карточной сети
до эмитента, хост последовательно снимает несколько рисков:
проверяет статус карты и срок её действия, смотрит доступность
средств или свободного кредитного лимита, сверяет лимитный
контур, оценивает признаки нетипичного поведения
(см. главу~\ref{ch:fraud}) и, для EMV-транзакций, верифицирует
криптограмму ARQC через HSM~\cite{emvco-book2,emvco-book3}.

Если все проверки пройдены, хост формирует ответ \texttt{00}
(Approved) и генерирует ARPC (Authorization Response
Cryptogram), которая возвращается карте для взаимной
аутентификации.

Граничный случай: CMS недоступен, но авторизационный хост
работает. Хост не может проверить статус карты (заблокирована?
перевыпущена?) и текущий лимит. Два варианта. Первый --
отклонить все транзакции: безопасно, но парализует
весь карточный портфель. Второй -- работать по кэшу:
хост хранит копию критических полей (статус, лимит,
floor limit) с TTL 5--15~мин и продолжает авторизовать
в пределах кэшированных лимитов. Транзакции, превышающие
кэшированный лимит, отклоняются; для остальных эмитент
принимает на себя риск одобрения по устаревшим данным.
После восстановления CMS хост синхронизирует фактические
остатки и фиксирует расхождения. На портфеле в миллионы
карт разница между 5~мин и 30~мин простоя CMS может
составить сотни одобренных транзакций по заблокированным
картам~\cite{emvco-book3}. Если хотя бы одна проверка не проходит,
хост возвращает соответствующий response code
(\texttt{51}: Insufficient Funds, \texttt{05}: Do Not Honor,
\texttt{14}: Invalid Card Number и так далее), и транзакция
отклоняется. Эмитент не ищет
один магический сигнал истины: он слой за слоем сужает область
доверия, пока не останется либо безопасное \enquote{да},
либо объяснимое \enquote{нет}.~\cite{emvco-book3}

\subsection*{Сценарии эмитента (issuer scripting)}

Решение эмитента не всегда заканчивается словами Approved
или Declined. Иногда вместе с ответом нужно изменить поведение
самой карты. Для этого EMV-стандарт предусматривает issuer
scripting -- механизм передачи команд карте через
авторизационный ответ. Script 71 выполняется до генерации
финальной криптограммы (TC или AAC), Script 72 после.~\cite{emvco-book3}
На практике issuer scripting используют, когда нужно обновить
офлайн PIN без визита в банк, сбросить офлайн-счётчики после
успешной онлайн-авторизации, заблокировать приложение прямо
во время контакта карты с терминалом или обновить параметры
управления риском, например лимиты на офлайн-транзакции.

Issuer scripting связывает авторизацию
и жизненный цикл карты в одну операцию, но его применение
ограничено: не все терминалы одинаково надёжно обрабатывают
скрипты, а при бесконтактной транзакции окно времени слишком
короткое, и сложное обновление может не успеть
завершиться (см. главу~\ref{ch:contactless}).

\subsection*{ACS: сервер аутентификации}

В интернет-платежах эмитенту мало проверить карту, лимит
и криптограмму. Нужно ещё понять, действительно ли покупку
инициировал держатель. Эту задачу решает ACS, компонент
эмитента, отвечающий за аутентификацию 3-D Secure
(см. главу~\ref{ch:3ds}). Когда мерчант инициирует проверку
3-D Secure, запрос приходит в ACS, и тот решает, можно ли
пропустить операцию по сценарию frictionless, без участия
держателя, или нужен challenge, то есть одноразовый код, биометрия
или подтверждение в мобильном приложении.~\cite{emvco-3ds,emvco-3ds-whitepaper}

Решение ACS строится на Risk-Based Authentication (RBA):
учитываются устройство, IP-адрес, история прошлых транзакций,
сумма, профиль мерчанта и другие признаки контекста. Качество
этого решения напрямую влияет на конверсию. Слишком строгий ACS
отправляет на challenge операции, которые можно было бы пропустить
без лишнего трения, и мерчант теряет покупателей. Слишком мягкий
ACS, наоборот, пропускает мошенничество и увеличивает убытки
эмитента.~\cite{emvco-3ds-whitepaper}

\subsection*{Управление токенами}

После выпуска карты её реквизиты начинают жить на пластике,
в телефоне, браузере и у мерчанта
в виде сетевых токенов (DPAN) (см. главу~\ref{ch:tokenization}).
Эмитент участвует в этом жизненном цикле с самого начала.
Он получает от провайдера токенизации (TSP) запрос
на tokenization, принимает решение об одобрении после
процедуры ID\&V (identification and verification) держателя,
а затем управляет состоянием токена: активирует его,
приостанавливает или удаляет.~\cite{emvco-payment-tokenisation,visa-vts,nspk-tsp}

Токен живёт как отражение исходной карты.
Если держатель блокирует карту, эмитент обязан
приостановить и связанные с ней токены. Если карта перевыпущена,
нужно обновить привязку токенов к новому PAN. Иначе платёжный
опыт распадается: держатель видит новую карту в приложении банка,
а у кошелька или мерчанта продолжает жить старое представление
о ней.~\cite{emvco-payment-tokenisation,visa-vts}

\subsection*{Работа с диспутами}

Держатель карты с претензией обращается не в карточную сеть
и не к эквайеру, а к эмитенту -- своему банку: \enquote{я не совершал эту покупку}
или \enquote{товар не получен}. Эмитент проверяет обоснованность
жалобы и, если основания есть, инициирует chargeback через
карточную сеть, направляя его эквайеру
(см. главу~\ref{ch:disputes}).~\cite{visa-core-rules,mc-chargeback-guide,visa-vcr-merchants}

На практике этот контур начинается с приёма и классификации
обращения. Затем эмитент сопоставляет ситуацию с правилами сети:
есть ли подходящий reason code, не пропущены ли сроки, достаточно
ли документов и соблюдены ли обязательные условия для chargeback.
После этого эмитент формирует chargeback-сообщение, а если эквайер
присылает representment и подкрепляет его compelling evidence,
эмитент должен решить, принять ли это объяснение или идти дальше,
к pre-arbitration. Dispute management -- продолжение авторизации задним числом: авторизация отвечает на вопрос \enquote{можно ли пустить деньги},
диспут -- на вопрос \enquote{надо ли это решение
пересмотреть}.~\cite{visa-core-rules,mc-chargeback-guide,visa-vcr-merchants}

\section{Эквайер}\label{sec:ia-acquirer}

Эквайер (acquirer) видит ту же самую операцию с противоположной
стороны. Он подключает терминал или платёжный шлюз, доставляет
авторизационный запрос в карточную сеть, принимает расчёт
от эмитентов и перечисляет деньги мерчанту за вычетом комиссий.
Эквайерская архитектура покрывает приём платежа,
удержания, выплату мерчанту и правила расчёта.

\subsection*{Система управления терминалами (TMS)}

\definitionnote{TMS}{Система удалённого управления парком
  POS-терминалов: обновления, ключи, конфигурация терминала,
  мониторинг состояния и диагностика неисправностей.}

В офлайн-эквайринге первый физический слой -- терминал.
У крупного эквайера TMS управляет
десятками и сотнями тысяч устройств, и главная задача здесь:
постоянно поддерживать актуальную конфигурацию: обновлённые
BIN-таблицы, свежие ключи, корректные EMV-параметры.

Без удалённой инъекции ключей (Remote Key Injection, RKI)
эквайринг плохо масштабируется:
ручная церемония загрузки ключей на каждый терминал
быстро упирается в операционный потолок.
RKI снимает это ограничение -- терминалы разворачиваются и обновляются
без физического доступа.~\cite{visa-tpa-registration-program-faqs}

\subsection*{Управляющий контур}

Для интернет-эквайринга и омниканальных сценариев одних терминалов
недостаточно. Нужен ещё слой, в котором живут профили мерчантов,
конфигурация MID/TID, правила дескрипторов, маршрутизация
по процессорам, параметры 3-D Secure, лимиты на возвраты
и выплаты, а также связка с платёжным шлюзом из главы~\ref{ch:gateway}. Такой слой часто называют управляющим
контуром.

Управляющий контур определяет, какой MCC и какой дескриптор увидит
эмитент, через какой BIN эквайера пойдёт транзакция и в какой
контур мониторинга попадёт мерчант. Если такие решения
расползаются между таблицами, тикетами и ручными настройками
в кабинете процессора, эквайер теряет управляемость: одна
и та же команда мерчанта начинает по-разному выглядеть для сети,
для команды риска и для финансовой
сверки.~\cite{visa-merchant-data-standards-manual}

У крупных эквайеров в управляющем контуре сходятся
подключение мерчанта, маршрутизация между несколькими процессорами
и операционная политика. Здесь решается, кому дать прямой
маршрут на конкретного процессора, кому оставить один канал,
а кому включить усиленный holdback или ручную проверку выплат.

\subsection*{Подключение мерчанта}

Эквайер не может включить приём платежей по одной только заявке
с сайта. Прежде чем дать мерчанту доступ к карточной сети, нужно
понять, кто перед ним, какой бизнес он ведёт и какой риск несёт.
Эту задачу решает контур подключения мерчанта --
набор проверок и настроек, определяющих и допуск,
и экономику отношений.

\textbf{KYB (Know Your Business)}: проверка юридического лица:
регистрация, бенефициары, финансовое состояние, репутация.
Эквайер обязан убедиться, что мерчант не занимается запрещённой
деятельностью, не находится в санкционных списках и ведёт
легитимный бизнес. Глубина проверки зависит от уровня риска:
интернет-казино проверяют строже, чем продуктовый магазин.

\textbf{MCC-классификация}: присвоение кода категории
(Merchant Category Code, MCC), четырёхзначного кода,
определяющего тип бизнеса.
MCC не формальность: от него зависит ставка interchange,
правила 3-D Secure, лимиты chargeback и иногда сама допустимость
операции. Например, MCC 7995 охватывает betting, lottery
и casino gaming, поэтому для таких операций сети и эквайеры
обычно вводят отдельные ограничения и усиленный
риск-контроль.~\cite{visa-merchant-data-standards-manual,visa-payment-facilitator-marketplace-risk-guide}

\practicenote{%
Неправильный MCC относится к самым частым ошибкам при подключении,
и последствия у неё вполне материальные. Мерчант может платить
неверный interchange, транзакции начинают некорректно проходить
антифрод-фильтры, а при аудите карточной сети ошибка приводит
к штрафу и требованию рекатегоризации. Выбор MCC требует той же внимательности, что и проверка юридических
документов.
}

\textbf{Underwriting}: оценка финансового риска мерчанта.
Эквайер смотрит, какой объём транзакций ожидается, каков средний
чек, насколько вероятны chargeback и какой кредитный лимит
можно предоставить до расчёта. Underwriting задаёт,
на каких условиях мерчант будет жить в системе: с какой ставкой
комиссии, с каким holdback, нужен ли rolling reserve и по какому
графику будут идти
выплаты.~\cite{visa-payment-facilitator-marketplace-risk-guide}
\subsection*{Мониторинг мерчанта}

После подключения риск не исчезает, и эквайер обязан
непрерывно наблюдать за мерчантом в рабочем режиме.
Мониторинг опирается на несколько групп сигналов:

\begin{itemize}
  \item \textbf{Chargeback ratio}: отношение числа chargeback
    к числу транзакций. Карточные сети устанавливают пороги
    и compliance-программы (для Visa сейчас это единый
      VAMP, объединивший мониторинг фрода и споров, а у Mastercard Excessive Chargeback
    Program, ECP, и Excessive Fraud Merchant, EFM)
    (см. главу~\ref{ch:disputes});
    их превышение ведёт к штрафам, усиленному мониторингу
    и в крайнем случае к отключению
    мерчанта.~\cite{visa-vamp-intro,mc-compliance-programs}
  \item \textbf{Fraud rate}: доля мошеннических транзакций.
  \item \textbf{Refund ratio}: аномально высокий процент
    возвратов, который может указывать либо на проблемы
    с товаром, либо на мошенническую схему
    (triangulation fraud).
  \item \textbf{Velocity anomalies}: резкий рост объёма
    транзакций, изменение среднего чека или появление
    нехарактерных BIN-диапазонов.
\end{itemize}

Хороший мерчант редко
становится проблемным за один день, но почти всегда заранее
оставляет следы в метриках: внезапный рост объёма, смещение
географии, неестественные возвраты, ухудшение доли
одобрений.~\cite{visa-payment-facilitator-marketplace-risk-guide}

\importantnote{%
Ответственность за мерчанта лежит на эквайере. Если мерчант
оказался мошенником или нарушил правила карточной сети,
штрафы и финансовые потери несёт именно эквайер, а не сеть
и не эмитент. Поэтому мониторинг мерчанта не формальность,
а защита от прямых финансовых потерь: пропущенный сценарий bust-out, когда
мерчант сначала наращивает объём, а затем исчезает с деньгами,
может стоить эквайеру
миллионов.~\cite{visa-payment-facilitator-marketplace-risk-guide,visa-acceptance-entities}

}
\subsection*{Ответственность за расчёт}

Для мерчанта эквайринг заканчивается в момент поступления денег.
Эквайер собирает клиринг, удерживает комиссии, применяет rolling reserve и
перечисляет средства мерчанту.

Эквайер должен уметь прозрачно объяснить каждую выплату:
какая сумма удержана как комиссия, какая отправлена в резерв,
какая зачтена против возвратов и в какой валюте считается
обязательство. Кроме того, он несёт риск late presentment,
временных разрывов по chargeback и кассовых разрывов между
получением денег от сети и выплатой мерчанту. Именно
эквайер задаёт политику выплат: T+1, T+2, еженедельные выплаты,
отложенный расчёт для high-risk MCC -- правила, которые
для мерчанта нередко важнее самой доли одобрений.

Holdback и rolling reserve покрывают временной разрыв
между авторизацией и закрытием окна возможных
chargeback -- до закрытия окна может пройти несколько месяцев. Если мерчант за это время
прекратил деятельность или не может покрыть возвраты, потери падают
на эквайера. Rolling reserve -- буфер, из которого эквайер покрывает
эти потери. Размер резерва отражает оценку риска мерчанта: высокий
chargeback rate, сезонный бизнес или новый мерчант без истории
получают более высокий процент удержания.

\section{Payment Facilitator и модель подмерчантов}\label{sec:ia-payfac}

Классическая модель эквайринга хорошо работает до тех пор,
пока мерчантов немного и каждого можно подключать как отдельного
клиента банка. Но если платформа хочет быстро вывести в платёжный
контур сотни или тысячи мелких продавцов, как маркетплейс
или SaaS-платформа с большим числом контрагентов, банковская процедура становится слишком
медленной. Индивидуальное подключение каждого продавца через банк
занимает дни или недели и начинает тормозить саму бизнес-модель.

\definitionnote{PayFac}{Агент спонсирующего эквайера, который
  подключает подмерчантов от его имени, получает расчёт в их пользу и
  отвечает за соблюдение правил сети и регуляторных требований.}

В модели PayFac платформа регистрируется спонсирующим эквайером
в карточной сети как Payment Facilitator и получает право
подключать sponsored merchants по ускоренной схеме. У подмерчанта
может быть собственный идентификатор, но договорная и расчётная
рамка всё равно идёт через связку PayFac + его эквайер: эквайер
перечисляет средства PayFac, а тот уже рассчитывается
с подмерчантами. Это и даёт главный выигрыш: нового продавца
можно подключить за минуты, а не за недели. PayFac автоматизирует
KYB, присваивает MCC и начинает приём платежей, не дожидаясь
полного банковского
цикла.~\cite{visa-acceptance-entities,visa-payment-facilitator-marketplace-risk-guide}

Однако скорость покупается концентрацией ответственности.
Именно PayFac отвечает перед карточной сетью и эквайером
за своих подмерчантов. Если один из них оказался мошенником,
штрафы ложатся на PayFac. Если у подмерчанта выросла доля
chargeback, отвечает тоже PayFac. Поэтому PayFac вынужден
строить собственные процессы KYB, мониторинга мерчантов
и работы с диспутами, по строгости мало отличающиеся
от банковских.~\cite{visa-acceptance-entities,visa-payment-facilitator-marketplace-risk-guide}

Главный структурный риск PayFac-модели: подмерчанты находятся вне
прямой видимости карточной схемы. Схема видит PayFac как единый
субъект; что происходит внутри его портфеля, она не знает напрямую.
Если подмерчант нарушает правила, схема обнаруживает это только через
отклонения в агрегированных показателях PayFac. Ответственность за
обнаружение -- на PayFac, финансовые последствия -- тоже на нём.

\importantnote{%
PayFac и ISO (Independent Sales Organization): разные роли.
ISO привлекает и обслуживает мерчантов
для банка-эквайера, получая агентское вознаграждение без прямого
участия в расчёте.
PayFac стоит прямо в денежной цепочке:
эквайер перечисляет средства ему, а он затем рассчитывается
с подмерчантами. Эта разница определяет и масштаб
ответственности, и набор регуляторных
требований.~\cite{visa-tpa-registration-program-faqs,visa-acceptance-entities}

}
\subsection*{Когда выбирать PayFac}

Такая модель оправдана, когда платформе нужно быстро
подключать большое число небольших мерчантов.
Маркетплейсы и подписочные платформы выигрывают от того,
что могут сами
управлять сценарием подключения, удерживать свою комиссию
и контролировать выплаты. Там, где мерчантов мало и они крупные,
прямой эквайринг чаще оказывается проще и дешевле.

\textbf{PayFac} (Payment Facilitator): платформа регистрируется
в сети через спонсирующего эквайера, ведёт KYB подмерчантов,
несёт ответственность за chargeback и контролирует выплаты.
Подходит при >100 подмерчантов и потребности
в быстром онбординге.
\textbf{MoR} (Merchant of Record): платформа -- единственный
мерчант в глазах сети; подмерчанты --
внутренняя абстракция платформы. Проще регуляторно, но
вся ответственность (chargeback, налоговые обязательства,
refund) замыкается на платформу. Подходит для SaaS
с единым продуктом.
\textbf{Pass-through} (прямой эквайринг): каждый мерчант
заключает договор с эквайером напрямую; платформа только
маршрутизирует трафик. Минимальная ответственность
платформы, но долгий онбординг (2--4 недели на мерчанта)
и невозможность управлять выплатами~\cite{visa-acceptance-entities}.
Если мерчантов <50 и они крупные -- pass-through.
Если мерчантов 50--500 и продукт единообразный -- MoR.
Если мерчантов >500 и нужен контроль над выплатами -- PayFac.

Цена этого удобства:
\begin{itemize}
  \item повышенная ответственность за мошенничество и chargeback
    подмерчантов;
  \item необходимость регистрации через спонсирующего эквайера
    в карточных сетях;
  \item инвестиции в собственные процессы KYB и AML;
  \item в некоторых юрисдикциях возникает необходимость дополнительной
    лицензии или специального статуса.
\end{itemize}

PayFac несёт эти обязанности, потому что встроен
в модель спонсорства эквайера и действует
по правилам карточной сети.~\cite{visa-acceptance-entities,visa-payment-facilitator-marketplace-risk-guide}

\practicenote{%
Если вы строите платформу и выбираете между моделью PayFac
и прямым эквайрингом для каждого мерчанта, оцените, сколько мерчантов вы реально планируете
подключить в первый год. Если речь идёт о десятках,
прямой эквайринг обычно проще и дешевле. Если о сотнях
и тысячах, PayFac начинает окупаться за счёт скорости
подключения, но требует серьёзных инвестиций в риск,
соблюдение требований и операционные процессы.
Промежуточный вариант:
использовать готовый PayFac-as-a-Service, например
Stripe Connect или Adyen for Platforms, и тем самым купить
скорость вместе с чужой инфраструктурой.

}
